\documentclass[11pt]{article}

\usepackage[margin=1in]{geometry}
\usepackage{amsmath}
\usepackage{amssymb}
\usepackage{amsthm}
\usepackage{booktabs}
\usepackage[hidelinks]{hyperref}

\newtheorem{theorem}{Theorem}
\newtheorem{lemma}{Lemma}
\setcounter{lemma}{-1} % the note's single lemma is Lemma 0
\newtheorem{corollary}{Corollary}
\newtheorem*{proposition}{Proposition}
\newtheorem{propositionnum}{Proposition}

\theoremstyle{definition}
\newtheorem{definition}{Definition}

\theoremstyle{remark}
\newtheorem{remark}{Remark}
\newtheorem*{remarkunnum}{Remark}

\newcommand{\Z}{\mathbf{Z}}
\newcommand{\Q}{\mathbf{Q}}
\newcommand{\R}{\mathbf{R}}
\newcommand{\C}{\mathbf{C}}
\newcommand{\Fun}{\mathbf{F}_1}
\newcommand{\NS}{\mathrm{NS}}
\newcommand{\Hom}{\mathrm{Hom}}
\newcommand{\End}{\mathrm{End}}
\newcommand{\Corr}{\mathrm{Corr}}
\newcommand{\Pic}{\mathrm{Pic}}
\newcommand{\pr}{\mathrm{pr}}
\newcommand{\covol}{\mathrm{covol}}
\newcommand{\trdeg}{\mathrm{trdeg}}
\newcommand{\SpecZ}{\operatorname{Spec}\Z}

\title{Products of the per-prime Tate curves of absolute geometry carry no
correspondence calculus for the Weil explicit formula}
\author{Jay Tyagi}
\date{August 26, 2026}

\begin{document}

\maketitle

\begin{abstract}
Connes and Consani's June-2026 absolute geometry of $\SpecZ$ produces, for
every prime $p$, a complex elliptic curve: the Tate curve
$E_p = \C^\times/p^{\Z}$, a rectangular torus with period lattice
$\Z \oplus \Z\cdot i\log p/(2\pi)$. A natural proposal is to assemble a
correspondence calculus for the Weil explicit formula --- the long-missing
``square of $\SpecZ$'' --- from the $(p,q)$-graded family of classical abelian
surfaces $E_p \times E_q$, importing Castelnuovo--Severi positivity fiber by
fiber. We prove that the proposal is empty. For distinct primes,
$\Hom(E_p, E_q) \neq 0$ if and only if $\log p \cdot \log q \in 4\pi^2\Q$
(Theorem~1), a transcendence coincidence that is impossible under the
four-exponentials conjecture. Unconditionally, each prime has at most one
exceptional partner prime (Theorem~2), so among $E_2\times E_3$,
$E_2\times E_5$, $E_3\times E_5$ at most one surface carries a nonzero
correspondence group (Theorem~3). Off this conjecturally empty matching,
$\NS(E_p \times E_q)$ has rank $2$: there are no graph classes and no diagonal
class, every divisor class acts as zero on $H^1$, and the
Castelnuovo--Severi inequality is vacuous (Theorem~4). On the diagonal,
$\End(E_p) = \Z$ unconditionally --- the CM alternative
$(\log p)^2 \in 4\pi^2\Q$ is refuted by Gelfond--Schneider --- and the
complete intersection profile of the graphs $\Gamma_m$ consists of universal
integers independent of $p$ (Theorem~5). No assignment of correspondence
groups, composition data, and intersection pairings on the family realizes the
commissioned calculus (Theorem~6): three independent obstructions, each
individually fatal. The family-level kill is unconditional; the note carries a
single open transcendence rider, an instance of the four-exponentials
conjecture, consumed by none of the theorems.
\end{abstract}

\setcounter{section}{-1}

\section{Summary}
\label{sec:summary}

Connes--Consani's June-2026 absolute geometry of $\SpecZ$ \cite{CC26}
produces, for every prime $p$, a complex elliptic curve --- the Tate curve
$E_p = \C^\times/p^{\Z}$, a rectangular torus with period lattice
$\Lambda_p = \Z \oplus \Z\cdot\tau_p$, $\tau_p = i\log p/(2\pi)$. For the
first time, the absolute-geometry program yields per-prime genus-$1$ fibers
carrying \emph{classical} intersection theory. The seed proposal examined
here was to assemble the long-missing correspondence calculus for the Weil
explicit formula --- the ``square of $\SpecZ$'' --- from the $(p,q)$-graded
family of classical abelian surfaces $E_p \times E_q$, importing the
Castelnuovo--Severi positivity fiber-by-fiber instead of waiting for the
unbuilt characteristic-one square.

\textbf{This note proves that the proposal is empty --- not approximately,
not ``bandwidth data in geometric dress,'' but empty.} The results:

\begin{enumerate}
\item \textbf{Isogeny criterion (Theorem~\ref{thm:isogeny}, iff):} for
distinct primes $p \neq q$, $\Hom(E_p, E_q) \neq 0$ \textbf{iff}
$\log p \cdot \log q \in 4\pi^2\Q$ --- a transcendence coincidence expected
never to occur (open unconditionally; impossible under the four-exponentials
conjecture, \S\ref{sec:transcendence}). When nonzero, $\Hom$ has rank $1$
and is computed exactly.
\item \textbf{Two-partner rigidity (Theorem~\ref{thm:twopartner},
unconditional):} each prime has at most \textbf{one} ``exceptional partner''
prime with which it could share a nonzero $\Hom$ --- pure unique
factorization, no transcendence input. The exceptional pairs, if any exist
at all, form a partial matching on the primes.
\item \textbf{The $\{2,3,5\}$ coherence kill (Theorem~\ref{thm:coherence},
unconditional):} among the three surfaces $E_2\times E_3$, $E_2\times E_5$,
$E_3\times E_5$, at most one carries a nonzero correspondence group; no two
of them can simultaneously do so. Any composition-coherent correspondence
family over all pairs of primes is unconditionally impossible.
\item \textbf{N\'eron--Severi collapse (Theorem~\ref{thm:NS}):}
$\NS(E_p \times E_q) = \Z\xi_1 \oplus \Z\xi_2 \oplus \Hom(E_p, E_q)$
\cite[Ch.~5]{BL}; reproved here. Off the exceptional matching, $\NS$ has
rank $2$ (the two fiber classes): there are \textbf{no graph classes, no
diagonal class}, every divisor class acts as \textbf{zero} on $H^1$, and the
Castelnuovo--Severi/Hodge-index inequality --- the positivity the seed was
commissioned to harvest ``for free'' --- is \textbf{vacuous} (its primitive
domain is the zero group).
\item \textbf{Diagonal residue (Theorem~\ref{thm:diagonal}):} the surfaces
$E_p \times E_p$ do carry a correspondence calculus ---
$\End(E_p) = \Z$ \textbf{unconditionally}, the CM alternative
$(\log p)^2 \in 4\pi^2\Q$ being impossible by Gelfond--Schneider
(\S\ref{sec:transcendence}) --- but its complete intersection profile
$(\Gamma_m\cdot\Gamma_n) = (m-n)^2$, $(\Gamma_m\cdot\Delta) = (m-1)^2$
consists of \textbf{universal integers independent of $p$}. The explicit
formula's transcendentally weighted prime terms $\Lambda(n) f(\log n)/\sqrt{n}$
cannot arise from it except by hand-inserting the weights into the
test-family coefficients --- re-entry into the known failure mode ``no new
data coordinate.''
\item \textbf{No-go (Theorem~\ref{thm:nogo}):} no assignment of
correspondence groups, composition data, and intersection pairing on the
family $\{E_p \times E_q\}$ realizes the commissioned calculus for the Weil
explicit formula. Three independent obstructions --- no composition material
(unconditional), no pairing target, prime-blind diagonal data --- each
individually fatal.
\end{enumerate}

Of the two transcendence questions the construction raises, only one survives
as a caveat: the diagonal coincidence $(\log p)^2 \in 4\pi^2\Q$ is
\textbf{refuted unconditionally} (Gelfond--Schneider;
\S\ref{sec:transcendence}, Proposition), and the off-diagonal coincidence
$\log p\cdot\log q \in 4\pi^2\Q$ --- (T1), an instance of the
four-exponentials conjecture --- is logged as the note's single open rider
(\S\ref{sec:transcendence}) with numerical evidence; \textbf{the family-level
kill is unconditional and consumes neither.}

Scope (\S\ref{sec:scope}): this kills the \emph{substrate family}, not the
direction. The Connes--Consani characteristic-one square (route~A) and the
Deninger leg (route~C) are untouched; the note in fact sharpens why the
square problem cannot be outflanked through the new complex fibers. Prior art
(\S\ref{sec:priorart}): nothing anticipates either the construction or the
obstruction; seven adjacent traditions are cited and distinguished.

\section{Background: the curves and the proposal}
\label{sec:background}

\subsection{The missing square}
\label{sec:missingsquare}

The classical template for proving a Riemann hypothesis --- Weil's positivity
argument on the square $C \times C$ of a curve over $\mathbf{F}_q$, with the
Castelnuovo--Severi inequality supplied by Riemann--Roch and ampleness ---
needs a doubled object. Over $\Q$ that object (``the square of $\SpecZ$,'' a
surface-like host on which the Weil explicit formula becomes an
intersection-number computation) exists in no branch of mathematics. Connes'
own summary of the stall, in print since 2015 and still accurate:
\emph{``[W]hat is missing is an intersection theory and a Riemann-Roch
theorem on the square of the arithmetic site''} \cite[\S4.3.2]{CE15}. The
square does exist in characteristic one --- Frobenius correspondences as
congruences on the square of the arithmetic site \cite{CC16} --- but with no
intersection theory and no Riemann--Roch, a decade on.

\subsection{The per-prime Tate curves}
\label{sec:tatecurves}

The June-2026 paper \cite{CC26} (``On the Absolute Geometry of Spec $\Z$ and
the Fargues--Fontaine curve'') changes the available raw material. Evaluating
the points of $(\SpecZ)_{\Fun}$ over $\C$, the non-trivial points over a
prime $p$ form a torsor under the Weil group $W_\infty = \C^\times$;
quotienting by the discrete Frobenius symmetries $p^\Z \subset W_\infty$
produces the \textbf{complex Tate curve}
\[
E_p := M_p^\infty / p^\Z \cong \C^\times / p^\Z ,
\]
\emph{``which canonically carries the structure of an elliptic curve over
$C$''} \cite[eq.~(4) and \S4.1.2]{CC26}. Lifting through
$w \mapsto \exp(2\pi i w)$ gives $E_p \cong \C/\Lambda_p$ with
\[
\Lambda_p = \Z \oplus \Z\cdot\tau_p , \qquad \tau_p = i\,\frac{\log p}{2\pi} ,
\]
a rectangular torus defined over $\R$ (the paper's own statement). Complex
conjugation descends to a canonical real structure whose real locus is two
copies of the adelic periodic orbit $C_p = \R^\times_{+}/p^\Z$. The paper's
Geometric Decomposition theorem (stated as Theorem~4 in its introduction and
Theorem~5 in \S4.3 of the full text) gives a canonical global splitting
\[
E_p \cong C_p \times \widetilde{X}_\infty ,
\]
where $C_p$ (length $\log p$) carries the entire dependence on $p$ and the
phase circle $\widetilde{X}_\infty$ (length $2\pi$, an unramified double
cover of the real Fargues--Fontaine analogue $X_\infty \cong \mathbf{P}^1(\R)$)
is \textbf{totally independent of $p$}; the modular parameter $\tau_p$ is
exactly the ratio of the two invariant $1$-forms' periods. We use three facts
and only these: the lattice $\Lambda_p$; the reality/rectangularity; and the
$p$-independence of the phase factor.

\textbf{Lineage (first-appearance record):} the per-prime
\emph{elliptic-analogue} idea first appears in the Connes--Consani orbit in
January 2025 --- \emph{``each $C_p$, when equipped with the restriction of
the scaling topos's structure sheaf, becomes an analogue of an elliptic curve
within the framework of characteristic one''} \cite[introduction]{CC25} ---
a characteristic-one (tropical) analogue of a single curve per prime: no
complex structure, no lattice, no products. \cite{CC26} complexifies it
seventeen months later. Neither paper raises products of the $E_p$,
correspondences between different primes, or any isogeny question ---
confirmed against the full texts (no isogeny/$\Hom$/product-of-Tate-curves
content).

\subsection{The seed proposal being killed}
\label{sec:seed}

The proposal under examination required, in the function-field template, a
calculus with three components:

\begin{itemize}
\item \textbf{(R1) Doubled object $+$ correspondence monoid.} Correspondence
groups containing the diagonal $\Delta$, the fiber classes $\xi_1, \xi_2$,
and one Frobenius-type class $\Psi_p$ per prime; closed under composition and
transpose; multiplicativity $\Psi_m \circ \Psi_n = \Psi_{mn}$ with $\Psi_p$
primitive exactly at primes.
\item \textbf{(R2) The pairing.} An intersection pairing with
$(\Psi_f \cdot \Delta) = W(f)$ for test elements
$\Psi_f = \int f(\lambda)\, \Psi_\lambda \, d^*\lambda$ --- the Weil explicit
formula as an intersection number, prime terms, archimedean term, and pole
term all accounted.
\item \textbf{(R3) Positivity.} A Hodge-index/Castelnuovo--Severi inequality
on the host, derived from Riemann--Roch plus effectivity (a non-circular
derivation, with no zeta input, is re-derived in a companion note; see the
Provenance remark).
\end{itemize}

The seed proposed to realize (R1)--(R3) on the $(p,q)$-graded family of
classical abelian surfaces $E_p \times E_q$, with the multiplicativity axiom
realized as a descent/compatibility datum \emph{across} the family, and the
archimedean and pole terms hoped to live on the $\widetilde{X}_\infty$
factor. The proposal's own stated fear was that the family might only
express \emph{bilinear} data --- $\Lambda(p^a)\Lambda(q^b)$ cross-terms,
bandwidth-style data in geometric dress. Theorems~\ref{thm:isogeny}--\ref{thm:nogo}
below show the situation is worse: the family expresses \textbf{no}
prime-facing correspondence data at all.

For the structure of the target, recall the explicit formula in one standard
normalization (Weil 1952 \cite{We52}; the exact archimedean shape is
immaterial here): for suitable test functions $f$, with
$\hat{f}(s) = \int_{\R} f(x)\, e^{(s-1/2)x}\, dx$,
\[
W(f) = \hat{f}(0) + \hat{f}(1)
 - \sum_{n\geq 2} \frac{\Lambda(n)}{\sqrt{n}}\bigl(f(\log n) + f(-\log n)\bigr)
 - W_\infty(f) ,
\]
$W_\infty$ an explicit archimedean integral. Only three structural features
are load-bearing for this note:

\begin{itemize}
\item \textbf{(W-i)} the prime side is \textbf{linear} in $\Lambda(n)$, with
transcendental weights $f(\pm\log n)/\sqrt{n}$ that vary with $n$ and with
$f$;
\item \textbf{(W-ii)} there is an archimedean term;
\item \textbf{(W-iii)} there are pole terms.
\end{itemize}

\section{Notation and the lattice presentation}
\label{sec:notation}

Throughout, $p, q, \dots$ denote primes; set
\[
y_p := \frac{\log p}{2\pi} > 0 , \qquad
\tau_p := i\, y_p , \qquad
\Lambda_p := \Z \oplus \Z\cdot\tau_p , \qquad
E_p := \C/\Lambda_p .
\]

\begin{lemma}[lattice presentation; {\cite[\S4.1.2]{CC26}}]
\label{lem:lattice}
$\exp_2\colon \C \to \C^\times$, $w \mapsto e^{2\pi i w}$, is a surjective
homomorphism with kernel $\Z$, and
\[
\exp_2(w + \tau_p) = e^{2\pi i w}\cdot e^{2\pi i \cdot (i y_p)}
 = e^{2\pi i w}\cdot e^{-\log p} = p^{-1}\cdot \exp_2(w) .
\]
Since $p^\Z = (p^{-1})^\Z$ as subgroups of $\C^\times$, $\exp_2$ induces an
isomorphism $\C/\Lambda_p \cong \C^\times/p^\Z$. \qed
\end{lemma}

(Numerically confirmed to 55 digits --- check~A of the numerical suite; see
the Provenance remark. The sign of $\tau_p$ is immaterial:
$\Z \oplus \Z(-\tau_p) = \Lambda_p$.)

Standard facts used without further comment: every homomorphism of complex
tori $\C/\Lambda \to \C/\Lambda'$ is induced by a unique $\lambda \in \C$
with $\lambda\Lambda \subseteq \Lambda'$ (lift to the universal cover); every
holomorphic map between complex tori is a homomorphism followed by a
translation (rigidity); $\Hom(E, E')$ denotes the group of holomorphic
homomorphisms.

\section{Theorem 1: the isogeny criterion (iff)}
\label{sec:isogeny}

\begin{theorem}
\label{thm:isogeny}
Let $p \neq q$ be distinct primes. Then
\[
\Hom(E_p, E_q) \neq 0
\iff y_p\, y_q \in \Q
\iff \log p \cdot \log q \in 4\pi^2\Q .
\]
Moreover, when the condition holds and $y_p y_q = u/v$ in lowest terms
($u, v \in \Z_{>0}$), the full homomorphism group is
\[
\Hom(E_p, E_q) = \Z\cdot\lambda_0 , \qquad \lambda_0 = i\, v\, y_q ,
\]
of rank $1$, and the isogeny induced by $\lambda_0$ has degree $uv$.
\end{theorem}

\begin{proof}
A nonzero homomorphism is multiplication by some $\lambda \in \C^\times$ with
$\lambda\Lambda_p \subseteq \Lambda_q$. Two containment conditions:

\begin{enumerate}
\item $\lambda\cdot 1 \in \Lambda_q$: so
$\lambda = a + b\cdot\tau_q = a + i\, b\, y_q$ for some $a, b \in \Z$.
\item $\lambda\cdot\tau_p \in \Lambda_q$: compute
$\lambda\cdot\tau_p = (a + i b y_q)\cdot(i y_p) = -b\, y_p y_q + i\, a\, y_p$.
Writing this as $c + i\, d\, y_q$ with $c, d \in \Z$:
\begin{itemize}
\item real part: $-b\, y_p y_q = c \in \Z$;
\item imaginary part: $a\, y_p = d\, y_q$.
\end{itemize}
\end{enumerate}

\emph{Case $a \neq 0$.} Then $d \neq 0$ (as $y_p, y_q > 0$) and
$y_p/y_q = d/a$, i.e.\ $\log p/\log q \in \Q$, i.e.\
$a\cdot\log p = d\cdot\log q$ with $a, d$ nonzero integers, i.e.\
$p^a = q^d$ (after clearing signs) --- impossible for distinct primes by
unique factorization. So $a = 0$, hence $d = 0$.

\emph{Case $a = 0$.} Then $\lambda = i\, b\, y_q$ with $b \neq 0$ (else
$\lambda = 0$), and the surviving condition is $b\, y_p y_q \in \Z$, which
forces $y_p y_q \in \Q$. This proves the forward direction, and identifies
$\Hom = \{\, i\, b\, y_q : b \in \Z,\ b\, y_p y_q \in \Z \,\}$. With
$y_p y_q = u/v$ in lowest terms, $b\, u/v \in \Z \iff v \mid b$, so
$\Hom = \Z\cdot(i\, v\, y_q)$.

\emph{Converse.} If $y_p y_q = u/v$, then $\lambda_0 = i\, v\, y_q$ satisfies
$\lambda_0\cdot 1 = v\cdot\tau_q \in \Lambda_q$ and
$\lambda_0\cdot\tau_p = (i v y_q)(i y_p) = -v\, y_p y_q = -u \in \Z \subseteq
\Lambda_q$, so $\lambda_0\Lambda_p \subseteq \Lambda_q$ and
$\lambda_0 \neq 0$ --- a genuine isogeny. (Numerically confirmed on a
synthetic pair with $y\, y' = 3/7$ exactly --- check~B.)

\emph{Degree.} $\covol(\Lambda_p) = y_p$ and
$\covol(\lambda_0\Lambda_p) = |\lambda_0|^2\cdot y_p = v^2 y_q^2\cdot y_p$,
so $\deg = [\Lambda_q : \lambda_0\Lambda_p] = v^2 y_q^2 y_p / y_q
= v^2\cdot(y_p y_q) = v^2\cdot(u/v) = uv$.

Finally $y_p y_q \in \Q \iff \log p \log q = 4\pi^2\cdot y_p y_q \in 4\pi^2\Q$.
\end{proof}

\begin{remark}[bookkeeping reconciliation]
\label{rem:bookkeeping}
The criterion may be stated as $\log p\cdot\log q \in 4\pi^2\Q$ or derived in
the form $b\cdot\log p\cdot\log q = -4\pi^2 c$ and phrased as
$\log p\cdot\log q \in \pi^2\Q$. These are the \textbf{same condition}:
$4\pi^2\Q = \pi^2\Q$ as subsets of $\R$, because $4 \in \Q^\times$ --- a
rational multiple of $\pi^2$ is a rational multiple of $4\pi^2$ and
conversely. The normalization-free statement is $y_p\cdot y_q \in \Q$; we
display the constant as $4\pi^2$ because
$y_p y_q = \log p\cdot\log q/(4\pi^2)$ makes it the natural one.
\end{remark}

\begin{remark}[symmetry]
\label{rem:symmetry}
The criterion is symmetric in $p$ and $q$; hence
$\Hom(E_p, E_q) \neq 0 \iff \Hom(E_q, E_p) \neq 0$ (as the general theory
also guarantees via the dual isogeny). ``Exceptional partner'' is therefore
a symmetric relation.
\end{remark}

\section{Theorems 2 and 3: the unconditional family-level kill}
\label{sec:familykill}

\begin{theorem}[two-partner rigidity]
\label{thm:twopartner}
For every prime $p$ there is \textbf{at most one} prime $q \neq p$ with
$\Hom(E_p, E_q) \neq 0$. Consequently the set of ``exceptional pairs''
$\{p, q\}$ (pairs with nonzero $\Hom$), if nonempty, forms a \textbf{partial
matching} on the primes: no prime appears in two of them.
\end{theorem}

\begin{proof}
Suppose $q \neq q'$ are two such primes. By Theorem~\ref{thm:isogeny},
$y_p y_q \in \Q^\times$ and $y_p y_{q'} \in \Q^\times$. Dividing:
$y_q / y_{q'} = (y_p y_q)/(y_p y_{q'}) \in \Q$, i.e.\
$\log q / \log q' \in \Q$, i.e.\ $q^m = q'^n$ for some positive integers
$m, n$ --- impossible for distinct primes by unique factorization.
\end{proof}

Nothing here consumes any transcendence input: the proof is unique
factorization plus linear algebra over $\Q$ applied to the products
$y_p y_q$, and it is agnostic about whether any exceptional pair exists.

\begin{theorem}[the $\{2,3,5\}$ coherence kill]
\label{thm:coherence}
Among the three surfaces $E_2\times E_3$, $E_2\times E_5$, $E_3\times E_5$,
at most one carries a nonzero correspondence group; in particular \textbf{no
two of them --- a fortiori not all three --- can simultaneously carry
nonzero correspondence groups.} (By Theorem~\ref{thm:NS} below,
``correspondence group of $E_p\times E_q$'' $= \Hom(E_p, E_q)$.)
\end{theorem}

\begin{proof}
Any two of the three pairs $\{2,3\}, \{2,5\}, \{3,5\}$ share a prime, so two
simultaneously nonzero groups would give some prime two exceptional
partners, contradicting Theorem~\ref{thm:twopartner}. Directly: if
$y_2 y_3 = r_1 \in \Q^\times$ and $y_2 y_5 = r_2 \in \Q^\times$ then
$y_3/y_5 = r_1/r_2 \in \Q$, i.e.\ $\log 3/\log 5 \in \Q$, i.e.\
$3^m = 5^n$ --- impossible; the other two cases are identical.
\end{proof}

\begin{remark}[enlarging the family does not help]
\label{rem:enlarge}
Define $E_n := \C^\times/n^\Z$ for any integer $n \geq 2$ (modulus
$y_n = \log n/(2\pi)$) --- the natural home for composite-index members
$\Psi_{pq}$ of the commissioned monoid. The computation of
Theorem~\ref{thm:isogeny} goes through verbatim, with one new phenomenon:
the case $a \neq 0$ now requires $\log m/\log n \in \Q$, which for general
integers means $m$ and $n$ are \textbf{multiplicatively dependent}
($m = k^a$, $n = k^b$), and then $E_m$, $E_n$ are indeed isogenous (the
subgroups $k^{a\Z}, k^{b\Z}$ of $\C^\times$ are commensurable). So the
enlarged family admits exactly the isogenies of \textbf{one-generator
towers} $\{E_{k^a}\}_a$ --- and across towers, by the two-partner argument
applied to multiplicatively independent moduli, at most the same
conjecturally-empty exceptional structure as in
Theorem~\ref{thm:twopartner}, now a partial matching of \emph{towers} (two
partners $n, n'$ of a fixed $m$ are forced into a single tower, and a
cross-tower isogeny would need the same unproven rational-product
coincidence as Theorem~\ref{thm:isogeny}). This is a structural echo of a
parallel finding on the same commissioned axiom system: the printed
polarized-Frobenius axioms are satisfiable precisely by one-generator
(single-prime) monoids --- the characteristic-$p$ template is a one-prime
template. The geometry says the same thing: per-prime towers exist; the
cross-prime structure, which is where the arithmetic of $\zeta$ lives, is
exactly what is absent.
\end{remark}

\begin{remark}[what a failure of the transcendence rider would and would not
change]
\label{rem:failure}
If some pair of primes actually satisfied
$\log p\cdot\log q \in 4\pi^2\Q$ --- spectacular transcendence news ---
Theorem~\ref{thm:isogeny} would produce one isogeny of computable degree
$uv$. It would not revive the seed: Theorems~\ref{thm:twopartner}
and~\ref{thm:coherence} cap the total supply of cross-prime maps at a
partial matching, and the obstructions O2 and O3 of Theorem~\ref{thm:nogo}
are untouched.
\end{remark}

\section{Theorem 4: N\'eron--Severi collapse off the matching}
\label{sec:NS}

We first fix the classical dictionary. For elliptic curves $E, E'$ over
$\C$, the group of \emph{divisorial correspondences} from $E$ to $E'$ is
$\NS(E \times E')$ modulo the subgroup generated by the two fiber classes
$\xi_1 = [\{0\}\times E']$, $\xi_2 = [E\times\{0\}]$ (Weil's ``trivial
correspondences''):
\[
\Corr(E, E') := \NS(E \times E') / \langle \xi_1, \xi_2 \rangle .
\]

\begin{theorem}
\label{thm:NS}
Let $E, E'$ be elliptic curves over $\C$. Then there is a canonical
isomorphism
\[
\NS(E \times E') \cong \Z\cdot\xi_1 \oplus \Z\cdot\xi_2 \oplus \Hom(E, E') ,
\]
hence $\Corr(E, E') \cong \Hom(E, E')$ (\cite[Ch.~5]{BL}; self-contained
proof below). Consequently, for distinct primes $p \neq q$ with
$\Hom(E_p, E_q) = 0$ (i.e., every pair outside the --- conjecturally
empty --- exceptional matching of
Theorems~\ref{thm:isogeny}--\ref{thm:twopartner}):
\begin{itemize}
\item[\textup{(a)}] \textbf{Rank collapse.}
$\NS(E_p \times E_q) = \Z\xi_1 \oplus \Z\xi_2$ has rank $2$, with
intersection form $\xi_1^2 = \xi_2^2 = 0$, $\xi_1\cdot\xi_2 = 1$: every
divisor class is numerically $a\xi_1 + b\xi_2$, and $\Corr(E_p, E_q) = 0$.
\item[\textup{(b)}] \textbf{$H^1$-inertness.} Every divisor class on
$E_p \times E_q$ acts as \textbf{zero} on $H^1$ (in both directions,
$H^1(E_p) \to H^1(E_q)$ and the transpose). The correspondence calculus is
not merely small; as a calculus of operators it is the zero calculus.
\item[\textup{(c)}] \textbf{No graphs, no diagonal.} The only holomorphic
maps $E_p \to E_q$ are constants, so the only graph classes are fiber
classes. There is no diagonal: not as a curve (a ``diagonal'' in a product
of nonisomorphic curves means the graph of an isomorphism, and none exists),
and not even as a numerical class --- no $D \in \NS$ satisfies the
diagonal's numerical triple $(D\cdot\xi_1, D\cdot\xi_2, D^2) = (1, 1, 0)$,
since $D\cdot\xi_1 = D\cdot\xi_2 = 1$ forces $D = \xi_1 + \xi_2$ and then
$D^2 = 2 \neq 0$.
\item[\textup{(d)}] \textbf{Castelnuovo--Severi vacuous.} The
Hodge-index/Castelnuovo--Severi inequality on a product of curves is the
statement that the intersection form is negative semidefinite on the
\emph{primitive} part $\langle \xi_1, \xi_2 \rangle^\perp \subset \NS$. Here
that primitive part is $0$: if $D = a\xi_1 + b\xi_2$ satisfies
$D\cdot\xi_1 = b = 0$ and $D\cdot\xi_2 = a = 0$ then $D = 0$. Equivalently:
every class of bidegree $(d_1, d_2)$ satisfies $D^2 = 2 d_1 d_2$ ---
Castelnuovo--Severi holds \emph{with equality everywhere}, carrying zero
information. The positivity the seed advertised as ``free fiber-by-fiber''
is free because there is nothing left for it to constrain.
\end{itemize}
\end{theorem}

\begin{proof}[Proof of the decomposition]
All inputs are standard abelian-variety theory (\cite[Ch.~2]{BL} for the
Poincar\'e bundle, seesaw, and rigidity; Ch.~5 for the endomorphism context;
the statement itself is the classical correspondence dictionary). Write
$\hat{E}' = \Pic^0(E')$ for the dual, $P$ for the Poincar\'e bundle on
$\hat{E}' \times E'$, normalized so that
$P|_{\{\alpha\}\times E'} \cong \alpha$ for $\alpha \in \hat{E}'$ and
$P|_{\hat{E}'\times\{0\}}$ is trivial.

\emph{Step 1 (the homomorphism $\varphi_L$).} For $L \in \Pic(E\times E')$
put $L_x := L|_{\{x\}\times E'}$; $\deg L_x$ is constant in $x$. Normalize:
\[
\widetilde{L} := L \otimes \pr_1^*\bigl(L|_{E\times\{0\}}\bigr)^{-1}
 \otimes \pr_2^*(L_0)^{-1} ,
\]
so that $\widetilde{L}_x = L_x \otimes L_0^{-1} \in \Pic^0(E')$ for all $x$,
$\widetilde{L}_0$ is trivial, and $\widetilde{L}|_{E\times\{0\}}$ is
trivial. By the universal property of $(\hat{E}', P)$ there is a unique
morphism $\varphi_L\colon E \to \hat{E}'$ with
$(\varphi_L \times \mathrm{id}_{E'})^* P \cong \widetilde{L}$ (the a-priori
ambiguity by $\pr_1$-pullbacks is killed by restricting to $E\times\{0\}$,
where both sides are trivial). Since $\widetilde{L}_0$ is trivial,
$\varphi_L(0) = 0$; a pointed morphism of abelian varieties is a
homomorphism (rigidity). So $\varphi_L \in \Hom(E, \hat{E}')$.

\emph{Step 2 (kernel; factorization through $\NS$).} $L \mapsto \varphi_L$
is additive: normalization and restriction are additive in $L$, and
uniqueness plus the biadditivity of the Poincar\'e bundle
($((\mathrm{add})\times\mathrm{id})^* P \cong \pr_{13}^* P \otimes
\pr_{23}^* P$ on $\hat{E}'\times\hat{E}'\times E'$, \cite[Ch.~2]{BL}) give
$\varphi_{L\otimes L''} = \varphi_L + \varphi_{L''}$. If
$L \in \Pic^0(E\times E')$ then, since the dual of a product is the product
of the duals, $L = \pr_1^* M \otimes \pr_2^* M'$ with $M \in \Pic^0(E)$,
$M' \in \Pic^0(E')$; the normalization annihilates it, so $\varphi_L = 0$
and the map factors through $\NS(E\times E')$. Conversely if
$\varphi_L = 0$ then $\widetilde{L}_x$ is trivial for every $x$ and
$\widetilde{L}|_{E\times\{0\}}$ is trivial, so by the seesaw principle
$\widetilde{L}$ is trivial; hence
\[
L \cong \pr_1^*\bigl(L|_{E\times\{0\}}\bigr) \otimes \pr_2^*(L_0) ,
\]
whose N\'eron--Severi class is
$\deg\bigl(L|_{E\times\{0\}}\bigr)\cdot\xi_1 + \deg(L_0)\cdot\xi_2$ (the
class of $\pr_1^*$ of a degree-$d$ bundle is
$d\cdot[\{\mathrm{pt}\}\times E'] = d\cdot\xi_1$, and symmetrically). So the
kernel of $\NS \to \Hom(E, \hat{E}')$ is exactly $\Z\xi_1 \oplus \Z\xi_2$ (a
free rank-$2$ group: the intersection numbers against $\xi_2, \xi_1$ read
off the coefficients).

\emph{Step 3 (surjectivity and splitting).} For $\psi \in \Hom(E, \hat{E}')$
set $L^\psi := (\psi \times \mathrm{id}_{E'})^* P$. Then
$L^\psi_x = P|_{\{\psi(x)\}\times E'} = \psi(x) \in \Pic^0(E')$, $L^\psi_0$
is trivial, and $L^\psi|_{E\times\{0\}} = \psi^*\bigl(P|_{\hat{E}'\times\{0\}}\bigr)$
is trivial --- so $L^\psi$ is already normalized and
$\varphi_{L^\psi} = \psi$ by uniqueness. Additivity of
$\psi \mapsto [L^\psi]$ is the Poincar\'e biadditivity again. This splits
the sequence:
\[
\NS(E\times E') \cong \Z\xi_1 \oplus \Z\xi_2 \oplus \Hom(E, \hat{E}') .
\]

\emph{Step 4 (eliminating the dual).} For an elliptic curve,
$x \mapsto [\mathcal{O}_{E'}([x] - [0])]$ is a canonical isomorphism
$E' \cong \hat{E}'$ (the degree-one principal polarization; Abel--Jacobi),
so $\Hom(E, \hat{E}') \cong \Hom(E, E')$.
\end{proof}

\begin{proof}[Proof of \textup{(a)}--\textup{(d)}]
(a) is immediate from the decomposition with $\Hom = 0$, plus
$\xi_1^2 = \xi_2^2 = 0$ (disjoint parallel fibers) and
$\xi_1\cdot\xi_2 = 1$ (one transverse point).

(b) The action of a correspondence on cohomology,
$\alpha \mapsto \pr_{2*}([D] \cup \pr_1^*\alpha)$, depends only on the class
$[D] \in H^2(E\times E', \Z)$. K\"unneth:
$H^2 = (H^2\otimes H^0) \oplus (H^1\otimes H^1) \oplus (H^0\otimes H^2)$. A
class $\pr_1^*\beta$ ($\beta \in H^2(E)$) acts on $\alpha \in H^1(E)$ by
$\pr_{2*}(\pr_1^*(\beta \cup \alpha)) = 0$ since
$\beta \cup \alpha \in H^3(E) = 0$; a class $\pr_2^*\beta$ acts by
$\beta \cup \pr_{2*}(\pr_1^*\alpha) = 0$ since fiber integration drops the
degree of $\pr_1^*\alpha \in H^1$ below zero. Now $\xi_1 = \pr_1^*[\mathrm{pt}]$
and $\xi_2 = \pr_2^*[\mathrm{pt}]$, so with $\NS = \Z\xi_1 \oplus \Z\xi_2$
every \emph{algebraic} class in $H^2$ has vanishing
$H^1\otimes H^1$-component and acts as zero on $H^1$. The transpose
direction is identical using $\Hom(E_q, E_p) = 0$
(Remark~\ref{rem:symmetry}).

(c) Rigidity: every holomorphic map $E_p \to E_q$ is a homomorphism followed
by a translation; $\Hom = 0$ leaves the constants, whose graphs are
horizontal fibers $E_p \times \{c\} \equiv \xi_2$. The numerical computation
is displayed in the statement.

(d) Displayed in the statement; the classical Castelnuovo--Severi inequality
itself (for products of curves with actual correspondences) is re-proved
from Riemann--Roch $+$ ampleness, with no zeta input, in a companion note
(see the Provenance remark) --- the engine is a real theorem whose domain of
application is, on these surfaces, the zero group.
\end{proof}

\begin{remark}[what still exists on these surfaces]
\label{rem:stillexists}
Rank-$2$ $\NS$ does not mean the surfaces are curve-poor: $\xi_1 + \xi_2$ is
ample (Nakai: positive square, positive on every effective class), and
general members of $|3(\xi_1+\xi_2)|$ are irreducible curves dominating both
factors. The point of (a)--(b) is that every such curve is \emph{numerically}
a fiber combination and acts as zero on $H^1$: as a correspondence it is
invisible. The collapse is a collapse of the correspondence calculus, not of
the surface's geometry --- which is exactly what kills the seed, whose
entire value proposition was the calculus.
\end{remark}

\section{Theorem 5: the diagonal residue is prime-blind --- and the no-go}
\label{sec:diagonal}

The one place the family genuinely has correspondences is the diagonal line
$E_p \times E_p$.

\begin{theorem}[diagonal residue]
\label{thm:diagonal}
Fix a prime $p$.
\begin{enumerate}
\item $\End(E_p) = \Z$, \textbf{unconditionally}. (The lattice computation
below shows the only alternative would be $(\log p)^2 \in 4\pi^2\Q$ ---
which would make $\tau_p$ imaginary quadratic and $E_p$ a CM curve --- and
that coincidence is impossible: Gelfond--Schneider,
\S\ref{sec:transcendence}, Proposition. No $E_p$ has complex
multiplication.)
\item By (1), $\NS(E_p \times E_p) = \Z\xi_1 \oplus \Z\xi_2 \oplus \Z\cdot[\Delta]$
has rank $3$ and $\Corr(E_p, E_p) \cong \End(E_p) = \Z$, generated by the
diagonal. The available correspondences are (combinations of fiber classes
and) the graphs $\Gamma_m$ of multiplication by $m \in \Z$, with
$\Gamma_1 = \Delta$, and their complete intersection profile is:
\[
\Gamma_m\cdot\xi_1 = 1 , \quad
\Gamma_m\cdot\xi_2 = m^2 , \quad
(\Gamma_m\cdot\Gamma_n) = (m-n)^2 \ \ (m \neq n) , \quad
(\Gamma_m\cdot\Delta) = (m-1)^2
\]
(the last for all $m$: at $m = 1$ this is $\Delta^2 = 0$, by adjunction with
$K = 0$ and $e(E_p) = 0$).
\item \textbf{Every number in this profile is a universal integer,
independent of $p$.} More precisely: under the real-affine identification
$E_p \cong \R^2/\Z^2$ sending $(1, \tau_p)$ to the standard basis, the
subtori $\Gamma_m, \xi_1, \xi_2$ correspond to the \emph{same} subtori for
every $p$, and intersection numbers of divisor classes are topological (cup
products of integral classes). The map
$p \mapsto (\text{complete intersection profile of } E_p \times E_p)$ is
constant. The prime enters the surface only through its complex structure
$\tau_p$, and the complex structure enters the correspondence calculus only
through $\Hom$/$\End$ --- which is the same ring $\Z$ for every $p$
(part~1).
\end{enumerate}
\end{theorem}

\begin{proof}
(1) $\lambda\Lambda_p \subseteq \Lambda_p$ with $\lambda = a + i b y_p$
(from $\lambda\cdot 1 \in \Lambda_p$) requires
$\lambda\cdot\tau_p = -b y_p^2 + i a y_p \in \Lambda_p$, i.e.\
$-b y_p^2 \in \Z$ (and the $\tau_p$-coefficient $a \in \Z$, automatic). If
$b \neq 0$ then $y_p^2 \in \Q$, i.e.\ $(\log p)^2 \in 4\pi^2\Q$ ---
impossible by the \S\ref{sec:transcendence} Proposition
(Gelfond--Schneider; a $b \neq 0$ endomorphism would make
$\tau_p^2 = -y_p^2 \in \Q$, i.e.\ $\tau_p$ quadratic imaginary and $\End$ a
CM order, and no $E_p$ admits this). So $b = 0$ and $\lambda = a \in \Z$:
$\End(E_p) = \Z$ with no exceptional case.

(2) Theorem~\ref{thm:NS} with $E = E' = E_p$ gives
$\NS = \Z\xi_1 \oplus \Z\xi_2 \oplus \End(E_p)$, and the graph of
$\mathrm{id}$ is $\Delta$. Intersection numbers: $\Gamma_m$ meets a vertical
fiber $\{x\}\times E_p$ in the single transverse point $(x, mx)$, so
$\Gamma_m\cdot\xi_1 = 1$. It meets a horizontal fiber $E_p\times\{c\}$ in
the fiber $[m]^{-1}(c)$, which has $m^2$ points, all transverse ($[m]$ is
unramified), so $\Gamma_m\cdot\xi_2 = m^2$; here
$\#\ker([k]) = \#\bigl((\tfrac{1}{k}\Lambda_p)/\Lambda_p\bigr) = k^2$
(representatives $(i + j\cdot\tau_p)/k$, $0 \leq i, j < k$; check~D of the
numerical suite tabulates them for $k \leq 7$). For $m \neq n$,
$\Gamma_m \cap \Gamma_n = \{(x, mx) : (m-n)x = 0\} = \ker([m-n])$, with
$(m-n)^2$ points, each transverse (tangent directions $(1, m)$ and $(1, n)$
are independent), so $(\Gamma_m\cdot\Gamma_n) = (m-n)^2$; $n = 1$ gives
$(\Gamma_m\cdot\Delta) = (m-1)^2$ --- which is also the Lefschetz number
$1 - 2m + m^2$ of $[m]$ (trace on $H^0 - H^1 + H^2$), the integer Lefschetz
data of the family.

(3) The subtori $\Gamma_m$ are the images of $x \mapsto (x, mx)$, defined by
the \emph{integral} linear data $m$ alone; under the identification by the
basis $(1, \tau_p)$ they are the same integral subtori of $\R^4/\Z^4$ for
every $p$. Cup products of integral classes on a fixed topological manifold
are $p$-independent.
\end{proof}

\begin{theorem}[the no-go]
\label{thm:nogo}
No assignment of correspondence groups, composition data, and intersection
pairings on the family of classical surfaces $\{E_p \times E_q\}$ ($p, q$
prime; optionally enlarged by composite moduli $E_n$ per
Remark~\ref{rem:enlarge}) realizes the commissioned calculus
\textup{(R1)}--\textup{(R3)} for the Weil explicit formula. Specifically:
\begin{itemize}
\item[\textup{(O1)}] \textbf{The composition material does not exist
(unconditional at family level).} \textup{(R1)} needs one primitive $\Psi_p$
per prime with composites $\Psi_p \circ \Psi_q$ supported across the
family --- nonzero elements of cross-prime correspondence groups for
essentially all pairs. By Theorem~\ref{thm:NS},
$\Corr(E_p, E_q) \cong \Hom(E_p, E_q)$; by
Theorems~\ref{thm:twopartner}--\ref{thm:coherence} these groups vanish
outside at most a partial matching of exceptional pairs, and already over
$\{2,3,5\}$ at most one of the three groups is nonzero. The correspondence
category of the family (objects $E_p$, morphisms $\Corr$) is unconditionally
too disconnected to carry a multiplicative monoid indexed by all primes: the
multiplicativity axiom has no material to act on. (Whether even a single
exceptional pair exists is the open transcendence question (T1) of
\S\ref{sec:transcendence} --- unneeded for this conclusion.)
\item[\textup{(O2)}] \textbf{The pairing target does not parse off the
diagonal.} \textup{(R2)} pairs test correspondences against the diagonal. On
every off-matching surface $E_p \times E_q$ there is no diagonal and no
graph class --- not even a numerical class with the diagonal's profile
(Theorem~\ref{thm:NS}c) --- and every class acts as zero on $H^1$
(Theorem~\ref{thm:NS}b), so no trace-type functional is available; the
Castelnuovo--Severi inequality \textup{(R3)} is vacuous there
(Theorem~\ref{thm:NS}d). The positivity engine is a real theorem with an
empty domain.
\item[\textup{(O3)}] \textbf{The diagonal residue cannot host the prime
terms.} On the surfaces $E_p \times E_p$ the calculus exists but its
complete intersection profile is the universal integer data of
Theorem~\ref{thm:diagonal}, identical for every $p$. A test element
$\Psi_f = \sum_m c_m(p)\cdot\Gamma_m + a\cdot\xi_1 + b\cdot\xi_2$ has
\[
(\Psi_f \cdot \Delta) = \sum_m c_m(p)\cdot(m-1)^2 + a + b ,
\]
a fixed integer-coefficient linear functional of the coefficients. By
\textup{(W-i)} the explicit formula's prime side is linear in $\Lambda(n)$
with weights $f(\pm\log n)/\sqrt{n}$ --- transcendental, $n$-dependent,
$f$-dependent, and non-constant in $p$. Since the geometric data is constant
in $p$, \emph{any} dependence of $(\Psi_f\cdot\Delta)$ on $p$ must be
written into the coefficients $c_m(p)$ by hand --- i.e., the designer
inserts the answer key $\Lambda(p) f(\log p)/\sqrt{p}$ into the test family
and the geometry contributes nothing. This is precisely the classical
failure mode ``Weil positivity in disguise'' (no new data coordinate), the
failure mode the proposal's own terms promised to avoid: the prime contact
was to be ``the correspondence calculus itself or nothing.'' It is nothing.
The archimedean term \textup{(W-ii)} fares no better: the candidate home
named by the proposal --- the phase factor $\widetilde{X}_\infty$ --- is
$p$-independent \emph{by the source paper's own decomposition theorem}
(\cite[Thm.~5 and \S4.3]{CC26}: ``totally independent of $p$''), so it
cannot source prime-dependent weights; and the pole terms \textup{(W-iii)}
have no candidate home at all.
\end{itemize}
Each obstruction is individually fatal; they are logically independent
(\textup{O1} concerns the cross-prime category, \textup{O2} the off-diagonal
surfaces, \textup{O3} the diagonal ones). \qed
\end{theorem}

\begin{remarkunnum}[Rider: the escape routes forfeit the point]
The $p$-flow and $q$-flow are conjugate through the real-analytic maps
$\lambda e^{i\theta} \mapsto \lambda^{\log q/\log p}\cdot e^{i\theta}$ ---
but these are not holomorphic, and a smooth (or characteristic-one)
correspondence theory on the products gives up exactly the classical
$(1,1)$-class positivity (Castelnuovo--Severi via Riemann--Roch $+$
ampleness) that was the seed's entire selling point, landing back on the
missing characteristic-one square with its absent intersection theory
\cite[\S4.3.2]{CE15}. There is no version of the family that keeps both the
cross-prime maps and the classical positivity.
\end{remarkunnum}

\section{The transcendence status: (T2) refuted (Gelfond--Schneider), (T1)
the single open rider}
\label{sec:transcendence}

Two number-theoretic questions arise from Theorems~\ref{thm:isogeny}
and~\ref{thm:diagonal} and are logged here. One of them turns out not to be
open at all: classical transcendence theory refutes it, in this note's
favor.

\begin{itemize}
\item \textbf{(T1)} For distinct primes $p, q$: is
$\log p \cdot \log q \in 4\pi^2\Q$ possible? (Equivalently: can $E_p$ and
$E_q$ be isogenous?) --- \textbf{open unconditionally}; refuted under the
four-exponentials conjecture (below).
\item \textbf{(T2)} For a prime $p$: is $(\log p)^2 \in 4\pi^2\Q$ possible?
(Equivalently: can $E_p$ have complex multiplication?) --- \textbf{no;
refuted unconditionally:}
\end{itemize}

\begin{proposition}[no CM anywhere in the family]
For every prime $p$, $(\log p)^2 \notin 4\pi^2\Q$. Hence no $E_p$ has
complex multiplication, and $\End(E_p) = \Z$ for every prime,
unconditionally (Theorem~\ref{thm:diagonal}(1)).
\end{proposition}

\begin{proof}[Proof (Gelfond--Schneider, 1934)]
Suppose $(\log p)^2 = 4\pi^2\cdot r$ with $r \in \Q$. Both sides are
positive, so $r > 0$, and taking positive square roots,
$\log p = 2\pi\sqrt{r}$, i.e.\ $p = e^{2\pi\sqrt{r}}$. With the branch
$\log(-1) = i\pi$,
\[
p = e^{2\pi\sqrt{r}} = e^{(-2i\sqrt{r})\cdot(i\pi)} = (-1)^{-2i\sqrt{r}} .
\]
The base $-1$ is algebraic, $\neq 0, 1$; the exponent
$\beta = -2i\sqrt{r}$ is algebraic (a product of $i$ and the real algebraic
number $\sqrt{r}$), nonzero, and purely imaginary --- hence irrational,
whether or not $\sqrt{r}$ is itself rational. By the Gelfond--Schneider
theorem, $(-1)^\beta$ is transcendental. But it equals the integer $p$ ---
contradiction.
\end{proof}

(Equivalently, in the linear-forms language: (T2) says
$\log p = (-2i\sqrt{r})\cdot(i\pi)$ --- a $\overline{\Q}$-linear relation
between the two $\Q$-linearly independent logarithms $\log p$ and
$\log(-1) = i\pi$, one real and one imaginary; their ratio would be the
algebraic irrational $-2i\sqrt{r}$, exactly what Gelfond--Schneider forbids.
The point is that this \emph{quadratic} relation among $\log p, \pi$
\textbf{factors} into linear-forms territory because $\sqrt{r}$ is
algebraic. It is the classical family ``$e^{\pi\sqrt{d}}$ is
transcendental'' --- the circle of Gelfond's constant $e^\pi$ and the
Ramanujan constant --- specialized to $d = 4r$.)

\paragraph{Status of (T1).}
Known unconditionally: $\log p/\log q \notin \Q$ (unique factorization ---
this is all Theorems~\ref{thm:twopartner}, \ref{thm:coherence} consume), and
the Proposition above on the diagonal. The genuinely quadratic
\emph{cross-prime} relation does not factor the same way --- it involves two
independent logarithms --- and no proven theorem currently excludes it. It
is, however, an \textbf{instance of the four-exponentials conjecture}
\cite{FourExp}, not merely adjacent to it:

\textbf{(4EC $\Longrightarrow$ $\neg$(T1)).} Suppose
$\log p\cdot\log q = 4\pi^2\cdot u/v$ ($u/v \in \Q$, necessarily positive).
Take $x_1 = 2\pi i$, $x_2 = \log p$ --- $\Q$-linearly independent (compare
real and imaginary parts) --- and $y_1 = 1$, $y_2 = \log q/(2\pi i)$ ---
$\Q$-linearly independent ($y_2$ is purely imaginary and nonzero, hence
irrational). The four exponentials are $e^{x_1 y_1} = 1$,
$e^{x_1 y_2} = q$, $e^{x_2 y_1} = p$, and
$e^{x_2 y_2} = e^{\log p\cdot\log q/(2\pi i)} = e^{-2\pi i\cdot u/v}$, a
root of unity. All four are algebraic --- contradicting the
four-exponentials conjecture. \qed

The \textbf{proven} six-exponentials theorem does not obviously reach (T1),
for an amusing reason: it needs a third row $x_3$ with both $e^{x_3 y_1}$,
$e^{x_3 y_2}$ algebraic, and --- taking $x_3 = \log \ell$ for a third prime,
the natural candidate --- the algebraicity of $e^{x_3 y_2}$ would itself
require a \emph{second} rational-product coincidence
$\log \ell\cdot\log q \in 4\pi^2\Q$, which this note's own
Theorem~\ref{thm:twopartner} (two-partner rigidity) forbids alongside the
first. The rigidity that makes the no-go unconditional is the same rigidity
that keeps its one open rider out of the proven theorem's reach.

\paragraph{Conditional resolution (Schanuel; retained as the stronger
statement).}
Schanuel's conjecture implies the full algebraic independence of
$\log p, \log q, i\pi$ --- hence $\neg$(T1), and it re-proves the
Proposition. Take $x_1 = \log p$, $x_2 = \log q$, $x_3 = i\pi$. These are
$\Q$-linearly independent: in $a\cdot\log p + b\cdot\log q + c\cdot i\pi = 0$
the imaginary part forces $c = 0$, and then unique factorization forces
$a = b = 0$. Schanuel then gives
\[
\trdeg_{\Q} \Q\bigl(\log p, \log q, i\pi, e^{x_1}, e^{x_2}, e^{x_3}\bigr)
= \trdeg_{\Q} \Q\bigl(\log p, \log q, i\pi, p, q, -1\bigr) \geq 3 ,
\]
and since $p, q, -1$ are algebraic, $\log p, \log q, i\pi$ are algebraically
independent over $\Q$. In particular
$\log p\cdot\log q + r\cdot(i\pi)^2 \neq 0$ for every rational $r$ --- i.e.\
$\log p\cdot\log q \notin \pi^2\Q = 4\pi^2\Q$. \qed

\paragraph{The status ladder for (T1).}
Open unconditionally $\subset$ settled by the four-exponentials conjecture
(as an instance) $\subset$ settled by Schanuel (with full algebraic
independence). Nothing in Theorems~\ref{thm:twopartner}--\ref{thm:nogo}
consumes any rung of it: the family-level kill needs only unique
factorization, which is exactly why it was engineered to be independent of
(T1) --- and (T2) is no longer a rider at all.

\paragraph{Numerical evidence (honesty check).}
Rows 1--3 of Table~\ref{tab:numerics} are evidence on the open question
(T1); rows 4--5 concern (T2), which the Proposition settles, and are
retained not as evidence but as \textbf{detector controls confirming a
theorem} --- the instruments correctly find no rational value where
Gelfond--Schneider says none exists. Working precision is 400 digits.

\begin{table}[ht]
\centering
\footnotesize
\begin{tabular}{@{}lllll@{}}
\toprule
target & value (first 37 digits) & \begin{tabular}[t]{@{}l@{}}PSLQ\\ (coeffs $\leq 10^{30}$)\end{tabular} & \begin{tabular}[t]{@{}l@{}}CF terms to\\ $q_N > 10^{100}$\end{tabular} & \begin{tabular}[t]{@{}l@{}}max partial\\ quotient\end{tabular} \\
\midrule
$\log 2\cdot\log 3/(4\pi^2)$ & $0.0192890205998215679007346039996508659\ldots$ & no relation & 201 & 4591 \\
$\log 2\cdot\log 5/(4\pi^2)$ & $0.0282579044193212256166934388383576738\ldots$ & no relation & 201 & 1598 \\
$\log 3\cdot\log 5/(4\pi^2)$ & $0.0447877188535867805034209039422156347\ldots$ & no relation & 179 & 284 \\
$(\log 2)^2/(4\pi^2)$ & $0.0121700170136801879558172286287261109\ldots$ & no relation & 160 & 226 \\
$(\log 3)^2/(4\pi^2)$ & $0.0305723743263550882536179008172733061\ldots$ & no relation & 191 & 1521 \\
\bottomrule
\end{tabular}
\caption{PSLQ and continued-fraction scans of the five targets at 400-digit
working precision. The values are reproduced to their full recorded
precision; the generating suite is named in the Provenance remark.}
\label{tab:numerics}
\end{table}

PSLQ finds no integer relation $m\cdot x + n = 0$ with coefficients up to
$10^{30}$ for any target; the continued-fraction expansions, scanned until
the convergent denominator exceeds $10^{100}$, do not terminate and show no
anomalously large partial quotient --- so \textbf{any rational value of any
target has denominator exceeding $10^{100}$} (if $x = a/b$ with
$b \leq q_N$, its continued fraction terminates within the scanned,
precision-guaranteed range). Detector control in the positive direction: the
same PSLQ instance recovers $\log 8 = 3\cdot\log 2$ instantly (check~E).
Status of this evidence: for (T1) it is evidence, not proof --- and nothing
in Theorems~\ref{thm:twopartner}--\ref{thm:nogo} consumes it; for (T2) the
proof is the Proposition, and the numerics merely confirm the detectors
work. In the exceptional-coincidence direction, Remark~\ref{rem:failure}
records that even a true (T1) coincidence would create only a single isogeny
inside the matching bound and would leave obstructions O2 and O3 intact.

\section{Scope: what dies, what does not}
\label{sec:scope}

\paragraph{What dies (this note is the publishable record).}
\begin{itemize}
\item The Tate-curve assembly seed, as commissioned: struck. The
$(p,q)$-graded family cannot carry (R1) (no composition material ---
unconditional), (R2) (no pairing target off-diagonal; prime-blind data
on-diagonal), or a non-vacuous (R3).
\item The planned ``first empirical data on the seed'' sandbox over
$\{E_2\times E_3, E_2\times E_5, E_3\times E_5\}$: ill-posed twice over ---
Theorem~\ref{thm:coherence} shows no coherent correspondence assignment over
the three surfaces exists to be measured (unconditionally, at least two of
the three carry no correspondence class at all), and off the conjecturally
empty exceptional matching the surfaces contain no graph classes and no
diagonal to enumerate against (Theorem~\ref{thm:NS}).
\item The thesis that the family furnishes ``a genuinely new substrate
option the literature has never raised'': the option is genuinely new
(\S\ref{sec:priorart}) and genuinely empty.
\end{itemize}

\paragraph{What does not die.}
\begin{itemize}
\item \textbf{Route A (the Connes--Consani characteristic-one square).}
Untouched --- and complemented: the stall there is a \emph{missing}
intersection theory on an \emph{existing} square \cite{CC16}\cite[\S4.3.2]{CE15},
whereas this note certifies that the classical-complex escape hatch out of
characteristic one --- the one the June-2026 curves seemed to open --- is
empty. The two results bracket the problem from opposite sides: the char-one
square has correspondences (Frobenius congruences) but no positivity
calculus; the complex surfaces have the full classical positivity calculus
but no correspondences. What the explicit formula needs is \emph{both at
once}, and neither habitat supplies the pair.
\item \textbf{Route C (the Deninger leg).} Untouched (a conormal-current
transplant to the Witt space, pursued separately).
\item \textbf{Hodge index over adelic curves.} Untouched, instrument-track;
this note if anything sharpens its honest framing (the positivity engine is
real and RH-free, but an engine is only as good as its substrate).
\item \textbf{The curves $E_p$ themselves.} Nothing here diminishes
\cite{CC26}'s construction; the no-go concerns only the pairwise
\emph{products} as a correspondence substrate --- a use the paper itself
never proposes.
\item \textbf{Pair-indexed data as such.} The Kurokawa
absolute-tensor-product line (\S\ref{sec:priorart}, item 5) proves the
\emph{analytic} ``square of $\SpecZ$'' exists at the zeta-function level,
with genuine Euler products over pairs of primes. Jointly with this note,
the failure is localized precisely: pair-of-primes data is realizable
analytically; what is empty is the proposed \emph{classical-geometric host}
built from per-prime fibers.
\end{itemize}

\paragraph{Taxonomy.}
The result binds all per-prime-fiber substrate designs, with the following
executable test for any future per-prime-fiber assembly proposal: (i)
compute $\Hom$ between two named fibers by the one-page lattice criterion;
(ii) exhibit the diagonal/graph classes the calculus needs on the actual
$\NS$ group; (iii) check whether any intersection number can carry a
$\log p$ weight without hand-insertion. The process lesson is that this
gate --- exactly the one-page computation of \S\ref{sec:isogeny} --- costs
one page of lattice algebra and should run at proposal time.

\section{Prior art: citations and distinctions}
\label{sec:priorart}

A dedicated literature search (see the Provenance remark) found \textbf{no
source anticipating either the construction} (products of per-prime complex
Tate curves as an arithmetic-surface surrogate) \textbf{or the obstruction}
(the rigidity of Theorems~\ref{thm:isogeny}--\ref{thm:diagonal}) ---
unsurprising for the construction, since the $E_p$ are two months old with
no citing literature; notable for the obstruction, which is elementary
lattice arithmetic that could have been recorded about any family of tori
with moduli $i\cdot\log p/(2\pi)$ and was not. Seven adjacent traditions
must be cited and distinguished:

\begin{enumerate}
\item \textbf{Connes--Consani--Marcolli, ``The Weil Proof and the Geometry
of the Adeles Class Space'' \cite{CCM07}.} The
correspondence-calculus-for-the-explicit-formula tradition --- but on the
square of the \emph{one} global adele class space, with correspondences
built from graphs of the scaling action and analogs of degree/codegree
(\S7.1 there). Its \S10.2 contains the germ of our objects: the tori
$T_p \cong \R^\times_{+}/p^\Z$ as singular-fiber components in a
vanishing-cycles analogy --- one-dimensional \emph{real} tori, no complex
structure, no products, no intersection theory on products.
\textbf{Distinction:} one global noncommutative space squared versus a
$(p,q)$-family of classical surfaces; the no-go concerns only the latter.
(The search also checked the distinct CCM paper ``Fun with $\mathbf{F}_1$,''
arXiv:0806.2401 --- zero Tate/elliptic/product content.)
\item \textbf{Connes--Consani, ``Geometry of the arithmetic site''
\cite{CC16} $+$ the essay stall quote \cite[\S4.3.2]{CE15}.} The square
\emph{exists} in characteristic one, with Frobenius correspondences as
congruences --- and intersection theory/Riemann--Roch missing a decade on.
\textbf{Distinction:} complementary result, as in \S\ref{sec:scope}: this
note shows the classical-complex escape from that stall is empty, and the
escape's price (smooth/char-one correspondences) is the classical positivity
itself.
\item \textbf{The lineage inside the Connes--Consani orbit
[\cite{CC25} $\to$ \cite{CC26}].} First appearance of per-prime
elliptic-analogue language: January 2025, characteristic one, single curve
per prime \cite{CC25}; complex Tate curves: June 2026 \cite{CC26}, which
never raises products, isogenies, or correspondences between primes
(confirmed against the full text). \textbf{Distinction:} both construction
and obstruction of this note are additions to, not extractions from, that
lineage.
\item \textbf{Banaszak--Uetake \cite{BU1, BU2, BU3}.} Independent
operator-theoretic axiomatization of a Weil-style intersection calculus
(axioms INT1--INT3 on a Hilbert-space operator, modeled on Weil's
$C \times C$), whose \emph{existence} is equivalent to RH ($+$
semi-simplicity). \textbf{Distinction:} they posit the calculus abstractly
and prove existence $\Longleftrightarrow$ RH; this note kills the one
concrete classical-surface instantiation ever proposed from per-prime
geometric data. Opposite poles: abstract axioms with RH-equivalent
existence versus concrete surfaces with provably empty calculus.
\item \textbf{The Kurokawa absolute tensor product line
\cite{Ku92, Ak09, Ta20}.} Kurokawa's 1992 proposal and the proven Euler
products \emph{over pairs of primes} for $\zeta \otimes \zeta$ and
Dirichlet-$L$ tensor squares (Koyama--Kurokawa; Akatsuka;
Kurokawa--Wakayama; Tanaka) --- ``the square of $\SpecZ$'' at the
zeta-function level, zeros at sums $\rho_1 + \rho_2$. \textbf{Distinction
and sharpening:} the pair-indexed \emph{analytic} object exists and
factorizes; the proposed \emph{geometric} host from per-prime Tate curves is
empty. Jointly the two localize the failure of the square at the geometric
hosting step, not at pair-indexed data. (This line is also the published
calibration for what a pairwise-product family \emph{would} express if it
expressed anything: pair-supported data with summed zeros --- not $W(f)$.)
\item \textbf{Haran's arithmetical-surface program \cite{Ha91}.} The closest
non-Connes--Consani surrogate proposal:
$\SpecZ \times_{\mathrm{alg}} \SpecZ$ reduces to the diagonal (Haran: ``the
surface reduces to the diagonal!''), so Haran proposes \emph{defining}
intersection numbers of ``Frobenius divisors'' on the nonexistent surface by
analytic means, aiming at a two-dimensional Riemann--Roch. \textbf{Distinction:} no per-prime elliptic fibers, no
products of classical curves, no obstruction statement; Haran's program
postulates the calculus's output, ours tests a concrete candidate host and
finds it empty.
\item \textbf{Scholze $A_{\mathrm{inf}}$/Fargues--Fontaine adjacency.}
Nothing in that school produces cross-prime products of per-prime objects as
a correspondence substrate (shtuka legs and $X_S$-families are all
single-$p$), and no such no-go is published there. The one genuine contact
runs through \cite{CC26} itself: the phase factor $\widetilde{X}_\infty$ as
a \emph{real analogue} of the Fargues--Fontaine curve --- a
per-prime-to-archimedean bridge, not a cross-prime one.
\textbf{Distinction:} disjoint mechanism; cited for completeness of the
search record.
\end{enumerate}

\paragraph{Transcendence background (rider citations).}
The Gelfond--Schneider theorem (1934) for the \textbf{unconditional
refutation of (T2)} --- the diagonal coincidence factors into a
$\overline{\Q}$-linear relation between $\log p$ and $i\pi$, squarely inside
the classical theorem (\S\ref{sec:transcendence}, Proposition); Baker's
theorem for the general linear-forms background; the four-exponentials
conjecture \cite{FourExp}, of which \textbf{(T1) is an instance} (not merely
an adjacent open problem), with the proven six-exponentials theorem kept out
of reach by Theorem~\ref{thm:twopartner} itself
(\S\ref{sec:transcendence}); Schanuel's conjecture for the stronger
algebraic-independence resolution. The question (T1) appears to be
unrecorded in the literature in this form and is stated here as open; the
\emph{question} (T2) may be equally unrecorded, but its resolution is a
two-line specialization of a classical theorem and is claimed here only as
such.

\paragraph{Novelty statement.}
Both the construction (the $E_p \times E_q$ assembly proposal, born and
killed inside the program that produced this note) and the obstruction
(Theorems~\ref{thm:isogeny}--\ref{thm:nogo}) are claimed as new, with the
seven citation-and-distinguish obligations above discharged. The
obstruction's ingredients are elementary --- Birkenhake--Lange-level lattice
arithmetic plus unique factorization --- and the search confirmed the
assembled statements appear nowhere.

\section*{Provenance}

\begin{remarkunnum}[Provenance]
This note is the publishable record of a result obtained on August 26, 2026
inside a structured research program on the Riemann hypothesis (direction
``C3: geometric substrate''); the seed proposal killed here
(\S\ref{sec:seed}) originated inside the same program, where it is retained
verbatim as an immutable record. The kill was constructed independently by
two verifiers under a duplicate-blind protocol (identical mathematics, no
shared text), re-derived in full by an independent adjudicator --- whose
record documents the $4\pi^2\Q$ versus $\pi^2\Q$ reconciliation of
Remark~\ref{rem:bookkeeping} --- and re-proved from scratch for this note.
The numerical checks cited in the text (check~A: the 55-digit confirmation
of Lemma~0; check~B: the synthetic isogeny pair with $y y' = 3/7$; check~C:
the PSLQ/continued-fraction table of \S\ref{sec:transcendence}; check~D: the
kernel tabulations for $k \leq 7$; check~E: the positive control
$\log 8 = 3\log 2$) are implemented in the program files
\texttt{seed-no-go-checks.py} and \texttt{seed-no-go-checks.json} at
400-digit working precision; a referee suite reproduced them bit-for-bit and
re-implemented them independently with rigorous interval-arithmetic
continued fractions and instrument-independent constants, with a third
instrument as cross-check. An independent referee pass dated August 26, 2026
found one fatal defect, cutting in the note's own disfavor: the original
version of \S\ref{sec:transcendence} claimed that no proven theorem excludes
either coincidence (T1)/(T2), whereas (T2) has been excluded unconditionally
since 1934 by Gelfond--Schneider, and (T1) is settled already by the
four-exponentials conjecture, far short of Schanuel. The repair --- the
Proposition of \S\ref{sec:transcendence}, the unconditional form of
Theorem~\ref{thm:diagonal}(1), and the corrected status ladder --- is
incorporated in this revision, along with nine mechanical referee repairs;
the strengthening consequence is that $\End(E_p) = \Z$ for every prime with
no hedge, leaving exactly one open rider, (T1). A dedicated prior-art gate
completed August 26, 2026 (roughly ten search phrasings; three papers
fetched and read; five corpus items read in full-text; unreachable
resources logged) returned the verdict \emph{novel-with-citations},
discharged in \S\ref{sec:priorart}; Akatsuka's paper \cite{Ak09} is
characterized through Tanaka's introduction \cite{Ta20} and venue
confirmations, its full text being paywalled --- nothing load-bearing rests
on unread portions. Quotations from \cite{CC26}, \cite{CC25}, and
\cite{CE15} were verified against full-text copies on file. The
non-circular derivation of Castelnuovo--Severi from Riemann--Roch plus
ampleness, cited in \S\ref{sec:seed} and \S\ref{sec:NS}, is a companion
program note; the result itself is banked in the program's barrier taxonomy
(entry IV.10) together with the executable test of \S\ref{sec:scope} and
cross-links to neighboring entries.
\end{remarkunnum}

\begin{thebibliography}{CCM07}

\bibitem[CC26]{CC26} A. Connes, C. Consani, \emph{On the Absolute Geometry
of Spec $\Z$ and the Fargues--Fontaine curve}, arXiv:2606.06604 (June
2026).

\bibitem[CC25]{CC25} A. Connes, C. Consani, \emph{Knots, primes and class
field theory}, arXiv:2501.06560 (January 2025).

\bibitem[CC16]{CC16} A. Connes, C. Consani, \emph{Geometry of the
arithmetic site}, Adv. Math. (2016).

\bibitem[CE15]{CE15} A. Connes, \emph{An essay on the Riemann Hypothesis},
arXiv:1509.05576, \S4.3.2.

\bibitem[CCM07]{CCM07} A. Connes, C. Consani, M. Marcolli, \emph{The Weil
Proof and the Geometry of the Adeles Class Space}, arXiv:math/0703392
(Manin Festschrift, Birkh\"auser).

\bibitem[BU1]{BU1} G. Banaszak, Y. Uetake, \emph{Abstract intersection
theory and operators in Hilbert space}, arXiv:0908.2909; Commun. Number
Theory Phys. 5 (2011).

\bibitem[BU2]{BU2} G. Banaszak, Y. Uetake, \emph{Standard models of
abstract intersection theory for operators in Hilbert space},
arXiv:1210.3526.

\bibitem[BU3]{BU3} G. Banaszak, Y. Uetake, \emph{Abstract intersection
theory for operators in Hilbert space: geometric aspects}, Funct. Approx.
Comment. Math. 64.2 (2021) 251--265.

\bibitem[Ku92]{Ku92} N. Kurokawa's absolute tensor product proposal (1992)
and its proofs: Koyama--Kurokawa ($\zeta \otimes \zeta$);
Kurokawa--Wakayama; see also \cite{Ak09}, \cite{Ta20}.

\bibitem[Ak09]{Ak09} H. Akatsuka, \emph{The double Riemann zeta function},
Commun. Number Theory Phys. 3 (2009) 619--653.

\bibitem[Ta20]{Ta20} H. Tanaka, arXiv:2008.07752 (history of the Kurokawa
absolute tensor product line and the Dirichlet-$L$ case).

\bibitem[Ha91]{Ha91} S. Haran, \emph{Index theory, potential theory, and the
Riemann hypothesis}, in \emph{L-functions and Arithmetic} (Durham, 1989),
J.~Coates and M.~J.~Taylor (eds.), London Math. Soc. Lecture Note Ser.\ 153,
Cambridge Univ. Press, Cambridge, 1991, pp.~257--270 (MR1110396; DOI
10.1017/CBO9780511526053.010). The passage relied on here is p.~259:
$\SpecZ \times \SpecZ \cong \SpecZ$, ``the surface reduces to the
diagonal''; the analytic definition $\langle f,g \rangle := W(f * g^{*})$ of
intersection numbers of ``Frobenius divisors'' on that nonexistent surface;
and ``a two dimensional Riemann--Roch for $\SpecZ$ may very well exist''.
Quoted verbatim and expanded in the survey K.~Thas, \emph{A taste of Weil
theory in characteristic one}, in \emph{Absolute Arithmetic and
$\Fun$-Geometry}, K.~Thas (ed.), European Mathematical Society Publishing
House, Z\"urich, 2016, pp.~365--386; arXiv:1507.06480.

\bibitem[BL]{BL} C. Birkenhake, H. Lange, \emph{Complex Abelian Varieties},
Springer, Ch.~2 and Ch.~5. (Citation at chapter level; every fact consumed
from it is re-proved in \S\ref{sec:NS}, so nothing load-bearing rests on a
recalled page number.)

\bibitem[We52]{We52} A. Weil, \emph{Sur les ``formules explicites'' de la
th\'eorie des nombres premiers}, Comm. S\'em. Math. Univ. Lund (1952),
Tome Suppl\'ementaire, 252--265. Used only through the
normalization-hedged structural display of \S\ref{sec:seed}; small-support
calibration as in E. Bombieri (2000).

\bibitem[4EC]{FourExp} The four-exponentials conjecture; see S. Lang,
\emph{Introduction to Transcendental Numbers}, Addison-Wesley (1966), and
M. Waldschmidt, \emph{Diophantine Approximation on Linear Algebraic
Groups}, Springer (2000).

\end{thebibliography}

\end{document}
